\chapter{ニューロンの種類}
\label{sec:neurontypes}

\section{ブラックボックスとしてのニューロン}

860億個のニューロンが入った袋を渡され~\cite{Azevedo2009}、それを山に仕分けてほしいと頼まれたとしよう。どう仕分けるだろうか。

見知らぬICが詰まった袋を渡されたエンジニアなら、パッケージの形で仕分けたりはしない。部品の入力はいくつか、出力はいくつか、出力に何を出すか、を問うはずである。ニューロンも同じように、回路図上のブロックとして描いてみよう（図~\ref{fig:blackbox}）。

\begin{figure}[t]
\centering
\begin{tikzpicture}[>=Stealth,x=1cm,y=1cm]
% the box
\node[draw,very thick,minimum width=2.6cm,minimum height=2.6cm,fill=blue!6] (N) at (0,0) {};
\node at (0,0.55) {\small ニューロン};
\node at (0,-0.15) {$\displaystyle u=\sum_i w_i x_i$};
\node at (0,-0.8) {\small $y=\Theta(u-\theta)$};
% inputs
\foreach \k/\y/\lab in {1/1.05/{x_1},2/0.55/{x_2},3/0.05/{x_3}}{
  \draw[->,thick,blue!60!black] (-4.2,\y) -- (-1.3,\y);
  \node[left] at (-4.2,\y) {$\lab$};
  \node[above,blue!60!black] at (-2.1,\y-0.06) {\scriptsize $w_{\k}$};}
\node at (-4.45,-0.4) {$\vdots$};
\node at (-2.1,-0.4) {$\vdots$};
\draw[->,thick,blue!60!black] (-4.2,-1.05) -- (-1.3,-1.05);
\node[left] at (-4.2,-1.05) {$x_n$};
\node[above,blue!60!black] at (-2.1,-1.11) {\scriptsize $w_{n}$};
\draw[decorate,decoration={brace,amplitude=5pt},gray!60!black] (-4.9,-1.2) -- (-4.9,1.2);
\node[left,align=right,gray!40!black] at (-5.1,0) {\small 入力\\[-2pt]\small $n\sim10^3$--$10^5$};
% single output wire, then fan-out
\draw[very thick,red!70!black] (1.3,0) -- (3.2,0);
\foreach \x in {1.6,2.2,2.34,2.9}{\draw[thick,red!70!black] (\x,0) -- (\x,0.4);}
\node[above right,font=\tiny,gray!30!black,inner sep=1pt] at (1.42,0.45) {カチッ\dots\ カチッカチッ\dots};
\node[below,red!70!black,align=center] at (2.25,-0.05) {\scriptsize 出力は$y$ひとつ};
\fill[red!70!black] (3.2,0) circle (0.06);
\foreach \y in {1.3,0.65,0,-0.65,-1.3}{
  \draw[->,thick,red!70!black] (3.2,0) -- (5.0,\y);}
\draw[decorate,decoration={brace,amplitude=5pt},gray!60!black] (5.3,1.45) -- (5.3,-1.45);
\node[right,align=left,gray!40!black] at (5.5,0) {\small $y$のコピーを\\[-2pt]\small $10^3$--$10^4$個の\\[-2pt]\small 別のニューロンへ};
\end{tikzpicture}
\caption{エンジニアの目で見たニューロン。入力$x_i$が多数あり、それぞれに重み$w_i$がつく。内部では重みつき和$u$を計算し、閾値$\theta$と比較する（$\Theta$は階段関数で、引数が正なら$1$、そうでなければ$0$）。出力$y$はひとつで、スパイクの列であり、分岐して数千の受け手に複製される。}
\label{fig:blackbox}
\end{figure}

ブロックの出力は数値ではなく、短い電圧パルスの列である。「カチッ」、休止、「カチッカチッ」、長い休止。どのスパイクも高さは同じなので、情報はすべてスパイクが\emph{いつ}起きるかにある。ブロックの内部には、第一次の粗い近似として、加算器とコンパレータがある。入力が重みつきで足し合わされ、和が閾値を超えるとブロックはスパイクを出す。次章では、これを連続時間で正しく動くモデルに置き換える。

出力はひとつだが、配線は枝分かれし、どの枝も\emph{まったく同じ}信号のコピーを運ぶ。エレクトロニクスでいうファンアウト（fan-out）である。大脳皮質のニューロンではファンアウトは数千のオーダーである。ICの論理ゲートは、ふつう数個のゲートを駆動するだけである。

では、出力配線は受け手に\emph{何を}届けるのか。スパイクは各枝の末端まで走り、そこで化学信号に変わる。枝は細胞間の狭い隙間に特定の物質、\emph{神経伝達物質}をひとつまみ放出し、それが受け手の入力電流を変える。これで仕分けの基準が決まる。\emph{ニューロンの出力がどの神経伝達物質を放出するか}である。

\section{出力ひとつに信号の型ひとつ}

この仕分けが意味をもつのは、各ニューロンの出力の型がひとつに限られる場合だけである。実際にそうなのか。実験によれば、ほぼそうである。

\begin{keyblock}
\begin{proposition}[ニューロン1個に出力の型1種類]
\label{prop:dale}
ほとんどすべてのニューロンは主要な神経伝達物質を1種類だけもち、出力のすべての枝でそれを放出する。したがって、ニューロンの型はその主要な神経伝達物質で決まる~\cite{Strata1999}。
\end{proposition}
\end{keyblock}

ニューロンの型は、\emph{主要な神経伝達物質の英語名の最初の4文字}で表すことにする。たとえば、主要な神経伝達物質がグルタミン酸（glutamate）であるニューロンは\nt{GLUT}ニューロンである。すべての型を表~\ref{tab:types}にまとめた。

\begin{table}[t]
\centering
\footnotesize
\setlength{\tabcolsep}{3pt}
\renewcommand{\arraystretch}{1.15}
\begin{tabular}{>{\raggedright}p{1.0cm}>{\raggedright}p{2.05cm}>{\raggedright}p{1.75cm}>{\raggedright}p{1.6cm}>{\centering}p{0.7cm}>{\raggedright}p{1.55cm}>{\raggedright}p{1.8cm}>{\raggedright\arraybackslash}p{3.2cm}}
\toprule
型 & 主要な神経伝達物質 & バス & 速さ & 符号 & 脳内のニューロン数 & ニューロン1個の出力数 & 計算での役割 \\
\midrule
\nt{GLUT} & グルタミン酸 & アドレス & 速い・遅い & $+$ & $\sim8\cdot10^{10}$ & $\sim10^4$ & 主要な正の重み \\
\nt{GABA} & GABA & アドレス & 速い・遅い & $-$ & $\sim3\cdot10^{9}$ & $10^3$--$10^4$ & 主要な負の重み \\
\nt{GLYC} & グリシン & アドレス & 速い & $-$ & $10^6$--$10^7$ & $10^2$--$10^3$ & 脊髄と脳幹での負の重み。グルタミン酸受容体の一種の第二の鍵 \\
\nt{ACET} & アセチルコリン & 両方 & 速い・遅い & $\pm$ & $\sim10^6$ & 筋肉：$10$--$10^3$、脳：$\sim10^5$ & 筋肉への出力。脳内では学習率（仮説） \\
\nt{DOPA} & ドーパミン & ブロード\newline キャスト & 遅い & $\pm$ & $\sim5\cdot10^{5}$ & $10^5$--$10^6$ & 報酬予測誤差$\delta$ \\
\nt{SERO} & セロトニン & ブロード\newline キャスト & 遅い（受容体の一種は速い） & $\pm$ & $\sim3\cdot10^{5}$ & $\sim10^5$ & 計画の時間的視野（仮説） \\
\nt{HIST} & ヒスタミン & ブロード\newline キャスト & 遅い & $+$ & $\sim6\cdot10^{4}$ & 推定なし & 全体への「覚醒せよ」信号 \\
\nt{NORA} & ノルアドレナリン & ブロード\newline キャスト & 遅い & $\pm$ & $\sim5\cdot10^{4}$ & $\sim10^5$ & ゲイン（仮説） \\
\nt{ADRE} & アドレナリン & ブロード\newline キャスト & 遅い & $\pm$ & $\sim10^4$ & 推定なし & ニューロンは少数。脳内での役割は不明 \\
\nt{ASPA} & アスパラギン酸 & アドレス & 速い & $+$ & 推定なし & 推定なし & 弱い正の重み。役割には異論がある \\
\bottomrule
\end{tabular}
\caption{ニューロンの型。脳内のニューロン数の多い順。\emph{バス}：アドレスバスでは、神経伝達物質は特定の受け手への狭い隙間に放出される。ブロードキャストでは、少数の発信元ニューロンから広い領域へ届く（\ref{sec:buses}節）。\emph{速さ}：速いはイオンチャネル型受容体経由でミリ秒、遅いは代謝型受容体経由で数百ミリ秒以上。\emph{符号}は通常の符号である。最終的に符号を決めるのは受け手であり（\ref{sec:receiver}節）、$\pm$は両方の符号がありうることを示す。\emph{脳内のニューロン数}は、ヒトの脳全体に含まれるその型のニューロンの数のオーダーである。ニューロンの総数は約$8.6\cdot10^{10}$である~\cite{Azevedo2009}。\nt{GLUT}ニューロンのほとんど、約$7\cdot10^{10}$個は小脳の小さな入力ニューロンである。\nt{GLUT}と\nt{GABA}の数は、この計数と大脳皮質における\nt{GABA}ニューロンの割合から求めた~\cite{Hendry1987}。少数派の型で直接数えられているのは\nt{HIST}~\cite{Airaksinen1991}と、\nt{DOPA}ニューロンの二つの群のうち主要なほう~\cite{Pakkenberg1991}だけなので、\nt{DOPA}の値は下限である。\nt{SERO}は二つの部分的な計数の和である~\cite{Baker1990,Baker1991}。\nt{NORA}、\nt{GLYC}、\nt{ACET}、\nt{ADRE}の値は直接の計数がない粗いオーダー推定である。\emph{ニューロン1個の出力数}は、1個のニューロンがもつシナプス終末、つまり出力配線の枝上にある神経伝達物質の放出点の数のオーダーである（図~\ref{fig:blackbox}）。ブロードキャスト型では終末を直接数えた例はない。\nt{DOPA}の値は、1本の出力配線の枝分かれが標的領域のどれだけを覆うかから導いたものであり~\cite{Matsuda2009}、それ以外の型の推定はさらに粗い。\nt{ACET}には二つの場合がある。筋肉を制御するニューロンと、脳内のブロードキャスト型ニューロンである。}
\label{tab:types}
\end{table}

% the pie is a table-type float with a figure caption, so it can never float ahead of Table~\ref{tab:types}
\begin{table}[tb]
\makeatletter\def\@captype{figure}\makeatother
\centering
\definecolor{vizglut}{HTML}{2A78D6}
\definecolor{vizgaba}{HTML}{E34948}
\definecolor{vizrest}{HTML}{8A8986}
\begin{tikzpicture}[>=Stealth]
% pie: GLUT 96.4 %, GABA 3.6 % (13.0 degrees), the rest 0.006 % (a hairline at 0 degrees)
\fill[vizglut] (0,0) -- (13:1.9) arc (13:360:1.9) -- cycle;
\fill[vizgaba] (0,0) -- (0:1.9) arc (0:13:1.9) -- cycle;
\draw[white,line width=1pt] (0,0) -- (13:1.9);
\draw[vizrest,line width=1.2pt] (0,0) -- (0:1.9);
\node[white,align=center,font=\small] at (190:0.95) {\nt{GLUT}\\$96.4\,\%$};
\draw[gray!60!black] (6.5:1.75) -- (2.05,2.1);
\node[anchor=west,font=\small] at (2.05,2.1) {\nt{GABA} $3.6\,\%$};
% zoom box for the remaining types
\draw[densely dashed,vizrest] (1.9,0) -- (3.1,1.65);
\draw[densely dashed,vizrest] (1.9,0) -- (3.1,-2.2);
\draw[vizrest,rounded corners=3pt] (3.1,-2.2) rectangle (10.1,1.65);
\node[anchor=west,font=\scriptsize] at (3.2,1.4) {その他すべての合計：$\approx0.006\,\%$、対数目盛};
\foreach \code/\lg/\y/\val/\lx in {GLYC/6.477/1.0/{$10^6$--$10^7$}/4.05,ACET/6.0/0.6/{$\sim10^6$}/3.0,DOPA/5.699/0.2/{$\sim5\cdot10^5$}/2.699,SERO/5.477/-0.2/{$\sim3\cdot10^5$}/2.477,HIST/4.778/-0.6/{$\sim6\cdot10^4$}/1.778,NORA/4.699/-1.0/{$\sim5\cdot10^4$}/1.699,ADRE/4.0/-1.4/{$\sim10^4$}/1.0}{
  \node[anchor=east,font=\scriptsize] at (4.35,\y) {\nt{\code}};
  \fill[vizrest] (4.45,\y-0.13) rectangle ({4.45+\lg-3},\y+0.13);
  \node[anchor=west,font=\scriptsize] at ({4.5+\lx},\y) {\val};}
\draw[vizrest,thick] (7.45,1.0) -- (8.45,1.0);
\draw[vizrest,thick] (7.45,0.92) -- (7.45,1.08);
\draw[vizrest,thick] (8.45,0.92) -- (8.45,1.08);
\draw[gray!60!black] (4.45,-1.72) -- (8.45,-1.72);
\foreach \e in {3,4,5,6,7}{
  \draw[gray!60!black] ({4.45+\e-3},-1.72) -- ({4.45+\e-3},-1.66);
  \node[below,font=\scriptsize] at ({4.45+\e-3},-1.72) {$10^{\e}$};}
\end{tikzpicture}
\caption{表~\ref{tab:types}の推定による脳内のニューロンの型の割合。円はほぼ全体が\nt{GLUT}で占められ、\nt{GABA}は細い扇形、それ以外の型はすべて合わせても境界上の髪の毛一本にすぎない。右側はその髪の毛一本をニューロン数の対数目盛で拡大したものである。\nt{GLYC}の棒は$10^6$--$10^7$の範囲の中央まで伸び、線分が範囲全体を示す。\nt{ASPA}は推定値がないので図にない。}
\label{fig:typepie}
\end{table}

これらの山の大きさがどれほど不揃いかは図~\ref{fig:typepie}でわかる。ニューロン1000個のうち$964$個が\nt{GLUT}、$36$個が\nt{GABA}であり、それ以外の型はすべて合わせても10万個あたり約6個にすぎない。

GABAはもともと4文字の略称である。主要な神経伝達物質になることがほとんどない物質、つまり神経ペプチド、気体、脂質、プリンには独自のコードを与えない。これらについては付録~\ref{sec:othermediators}で扱う。命題~\ref{prop:dale}の「ほとんど」という語は重要である。多くのニューロンは、主要な神経伝達物質と一緒に補助的な物質を放出する~\cite{Yao2023}。二つ目の速い神経伝達物質を、しかも逆の符号のものを放出するニューロンもある。\nt{DOPA}ニューロンの一部はGABAも放出する~\cite{Tritsch2012}。さらに、同じ出力の異なる枝から異なる神経伝達物質が出るニューロンもある~\cite{Zhang2015}。これらはすべて測定されているので、「ニューロン1個に神経伝達物質1種類」は例外のある規則であって、法則ではない。それでも、最初の区分としては検証に耐える。発現している遺伝子によってニューロンを5000以上の群に分けたマウス脳の細胞アトラスでも、ニューロンの大きなクラスはやはり、まず主要な神経伝達物質で分かれる~\cite{Yao2023}。

この規則からは、脳を人工ニューラルネットワークとただちに区別する帰結が導かれる。人工ネットワークでは、同じニューロンがある受け手には正の重みを、別の受け手には負の重みをもちうる。重み$w_{ij}$は行列の中の単なる数だからである。脳では、作用の符号はおもに神経伝達物質で決まり、ニューロンの神経伝達物質は1種類である。したがって、ニューロン$j$がある受け手を興奮させるなら、ほかの受け手もすべて興奮させる。ニューロン$j$から出る重み行列の\emph{列}全体が同じ符号をもつ。
\begin{empheq}[box=\keybox]{equation}
\operatorname{sign} w_{ij} \;=\; s_j \quad\text{すべての受け手 } i \text{ について}, \qquad s_j\in\{+1,-1\}.
\label{eq:dalesign}
\end{empheq}
ニューロンは興奮性（$s_j=+1$）と抑制性（$s_j=-1$）に分かれ、これは結合ではなくニューロン自体の性質である。興奮性ニューロンからの負の結合がネットワークに必要なら、仲介役を通すしかない。興奮性ニューロンが抑制性ニューロンを興奮させ、それが標的を抑制する。つまり、1段分の遅延をともなう負の重みである。

神経伝達物質は100種類以上知られているが、脳の仕事のほとんどは6種類ほどが担っており、それらはまったく異なる二つのクラスに分かれる。

\section{二つのバス：アドレスバスとブロードキャスト}
\label{sec:buses}

コンピュータネットワークには、メッセージを届ける方法が二つある。一つ目は\emph{アドレス指定}、つまりユニキャスト（unicast）である。パケットは1人の送り手から特定の受け手へ速く届き、ほかの誰にも見えない。二つ目は\emph{ブロードキャスト}（broadcast）である。送り手は一つのメッセージをネットワークの全ノードに一度に送る。ニューロンは両方の方法を使い、神経伝達物質はまさにこの基準で分かれる（図~\ref{fig:phone_radio}）。

\begin{figure}[t]
\centering
\begin{tikzpicture}[>=Stealth,scale=0.95]
% ---- unicast: point-to-point wiring
\node[above] at (2.0,3.3) {\small \textbf{アドレスバス}：グルタミン酸とGABA};
\foreach \i in {0,...,4}{
  \fill[blue!60!black] (0.2+0.9*\i,2.5) circle (0.1);
  \fill[blue!60!black] (0.2+0.9*\i,0.6) circle (0.1);}
\foreach \a/\b in {0/0,0/1,1/1,1/3,2/2,3/2,3/4,4/4,2/0}{
  \draw[->,blue!60!black] (0.2+0.9*\a,2.38) -- (0.2+0.9*\b,0.72);}
\fill[red!70!black] (1.1,1.55) circle (0.09);
\pic[scale=1.2] at (1.75,1.9) {letter};
\draw[->,red!70!black,thick] (1.1,1.55) -- (0.28,0.7);
\draw[->,red!70!black,thick] (1.1,1.55) -- (1.95,0.72);
\node[align=center,below] at (2.0,0.2) {\small 数十億のニューロン、\\[-2pt]\small それぞれに固有の宛先、\\[-2pt]\small 時間：ミリ秒};
% ---- broadcast: one small source, global targets
\begin{scope}[xshift=7.2cm]
\node[above] at (1.8,3.3) {\small \textbf{ブロードキャスト}：ドーパミンなど};
\foreach \i in {0,...,8}{\fill[blue!60!black] (-0.2+0.5*\i,2.75) circle (0.08);}
\fill[orange!85!black] (1.8,0.35) ellipse (0.35 and 0.18);
\pic[scale=1.5] at (2.7,0.75) {mast};
\foreach \x in {-0.2,0.3,0.8,1.3,1.8,2.3,2.8,3.3,3.8}{
  \draw[->,orange!85!black] (1.8,0.5) .. controls (1.8,1.3) and (\x,1.6) .. (\x,2.62);}
\node[align=center,below] at (1.8,0.1) {\small 数千から数十万のニューロン、\\[-2pt]\small 一つのメッセージを全員に、\\[-2pt]\small 時間：数百ミリ秒以上};
\end{scope}
\end{tikzpicture}
\caption{二つの伝達方式。左はアドレスバスで、封筒に宛先を書く郵便に似ている。赤は1個のニューロンとその受け手2個を示す。右はブロードキャストで、ラジオ局に似ている。発信元のニューロンは小さな群であり、同じメッセージを全員が受け取る。}
\label{fig:phone_radio}
\end{figure}

\emph{アドレスバス}はグルタミン酸とGABAである。圧倒的多数のニューロンがこれらを放出する。各出力は特定の細胞につながり、神経伝達物質は接点の狭い隙間の中だけで、しかも速く作用する。1ミリ秒で受け手のイオンチャネルが開き、数ミリ秒ですべてが終わる。これは\emph{データ}バスであり、脳が計算しているものがここを通る。

\emph{ブロードキャスト}のバスはドーパミン、セロトニン、ノルアドレナリン、そして一部はアセチルコリンである。これらを放出するのはごく小さなニューロン群だが、その一つひとつの出力は信じられないほど広く枝分かれする。神経伝達物質は狭い隙間ではなく細胞間の空間に放出されることが多く、そこで広がって周囲のすべてに作用する。作用は遅い。チャネルを直接開くのではなく、受け手の内部で化学反応の連鎖を起こす。ここを通るのはデータではなく、ネットワークの広い領域に共通の\emph{制御パラメータ}であり、それがネットワークの動作モードを変える。ゲイン、学習率、「今のはよかった」という信号である。そのため、このような神経伝達物質は\emph{修飾物質}（モジュレータ）と呼ばれる。こうしたチャネルはそれぞれ脳全体でただ一つの数だと長く考えられてきた。近年の測定はこの描像を精密にした。チャネルは1本の配線ではなく、並列な線を束ねた小さなハーネスである。同じ神経伝達物質の発信元ニューロンはサブシステムに分かれ、それぞれに固有の受け手と固有のメッセージがあり、その場でどれだけ放出するかは局所の信号によっても調整される~\cite{Ren2018,Poe2020}。こうした線の数は数本から数十本であって数十億ではないので、チャネルはやはりブロードキャストである。

この章から絵を一枚だけ持ち帰るなら、これである。脳はデータをアドレス指定で伝える巨大なネットワークであり、その上に制御信号を運ぶブロードキャストチャネルがいくつか張られている。プログラマなら見慣れた構成に気づくだろう。計算と、少数の制御変数の組である。それぞれの制御変数はネットワークの広い領域に共通である。

\section{\nt{GLUT}：正の重み}

グルタミン酸は脳の主要な興奮性神経伝達物質である。ニューロンの出力がグルタミン酸を放出すると、受け手ではナトリウムを内側に通すチャネルが開き、受け手の膜電位が上がって、受け手自身がスパイクを出す確率が高まる。図~\ref{fig:blackbox}の言葉でいえば、これは\emph{正の}重み$w_i>0$をもつ入力である。

グルタミン酸を出すのは誰か。長距離でデータを伝えるほとんどすべてのニューロンである。大脳皮質のニューロンの4分の3から5分の4、そして脳の異なる領域をつなぐニューロンの大部分がそうである。脳のネットワークは、何よりもまずグルタミン酸結合のネットワークである。

工学的な細部を一つ。放出の後、隙間の中のグルタミン酸濃度は数千倍に跳ね上がって約1ミリモーラーに達し、時定数約1ミリ秒で減衰する~\cite{Clements1992}。グルタミン酸は隙間から拡散して出ていき、専用のポンプタンパク質が細胞内へくみ戻す。なぜか。通信線路で各シンボルの後に信号をゼロに戻すのと同じ理由である。神経伝達物質を取り除かなければ、次のスパイクは「使用中」の線路に到着し、数ミリ秒間隔のスパイクは融合してしまう。

\section{\nt{GABA}：負の重み}

重みがすべて正のネットワークを想像してほしい。平均して1個のスパイクが次のスパイクを1個より多く生むなら、活動は指数関数的に増え、数ステップでネットワーク全体に広がる。電子回路の技術者なら、ループゲインが1を超える正帰還ループだとわかるだろう。回路は飽和する。そうならないためには負の重みが必要である。その主な供給源が$\gamma$-アミノ酪酸、略してGABAである。

GABAは受け手で塩化物イオンを通すチャネルを開く。塩化物イオンの平衡電位は静止電位付近かそれより少し低いところにあるので、開いた塩素チャネルは膜を静止電位付近に保ち、電圧を閾値まで上げさせない。これは\emph{負の}重み$w_i<0$をもつ入力である。GABAニューロンは大脳皮質のニューロンの5分の1から4分の1を占める~\cite{Hendry1987}。その出力は短く、すぐ隣のニューロンを抑制する。

脳における抑制は、単なる引き算よりはるかに多くのことをする。ネットワーク全体を動作点に保つ負帰還として働くのである。正常に動作している大脳皮質では、ニューロンに入る興奮性電流と抑制性電流がほぼ正確に釣り合っていると考えられている。この状態は理論で予言され~\cite{vanVreeswijk1996}、生きた皮質での電流測定でも見えており~\cite{Haider2006}、ニューロンは常に閾値付近にいる。不安定点の近くで強い負帰還によって安定化した回路は、小さな信号に速く敏感に反応する。昔からある工学の手法である。

\section{\nt{DOPA}：誤差信号}

最初のブロードキャストチャネルはドーパミンである。ヒトのドーパミンニューロンは約50万個、およそ17万個に1個である。これはドーパミンニューロンの二つの群のうち主要なほうで数えられた値なので~\cite{Pakkenberg1991}、下限である。その出力は、行動を選択する領域へと広がっていく。

このチャネルは何を伝えるのか。答えは、報酬を受け取る動物でドーパミンニューロンのスパイクを記録する実験から得られる~\cite{Schultz1997}。動物にときどき予告なしにジュースを1滴与えると、ドーパミンニューロンは短いスパイクのバーストで応答する。次に、ジュースの前にランプを点けるようにする。いつもジュースの1秒前である。動物がこの関係を学習すると、ニューロンは\emph{ランプ}にバーストで応答するようになり、予定どおりの時刻に来たジュースそのものにはまったく応答しなくなる。そしてランプの後にジュースを与えないと、ジュースが来るはずだった瞬間にニューロンは\emph{黙り込み}、発火率は通常より下がる。

三つの場合すべてを説明する規則はこうである（図~\ref{fig:rpe}）。ニューロンが伝えるのは、どれだけよいかではなく、\emph{予想よりどれだけよいか悪いか}である。

\begin{figure}[t]
\centering
\begin{tikzpicture}[>=Stealth,x=3.4cm,y=1cm]
\foreach \y/\lab in {2.8/{予告なしのジュース},1.4/{ランプの後のジュース},0/{ランプだがジュースなし}}{
  \draw[gray!70!black] (0,\y) -- (3,\y);
  \draw[densely dotted,gray!60!black] (0,\y+0.3) -- (3,\y+0.3);
  \node[left,align=right,font=\small] at (-0.03,\y+0.35) {\lab};}
% row 1: juice without a cue -> burst at the juice
\draw[very thick,orange!85!black] plot[domain=0:3,samples=120] (\x,{2.8+0.3+0.75*exp(-((\x-2.08)/0.07)^2)});
\pic[scale=0.8] at (2,2.58) {juice};
% row 2: cue then juice -> burst at the cue, nothing at the juice
\draw[very thick,orange!85!black] plot[domain=0:3,samples=120] (\x,{1.4+0.3+0.75*exp(-((\x-1.08)/0.07)^2)});
\pic[scale=0.85] at (1,1.15) {lamp}; \pic[scale=0.85] at (2,1.15) {juice};
% row 3: cue, juice omitted -> burst at the cue, dip when the juice was due
\draw[very thick,orange!85!black] plot[domain=0:3,samples=120] (\x,{0.3+0.75*exp(-((\x-1.08)/0.07)^2)-0.28*exp(-((\x-2.1)/0.12)^2)});
\pic[scale=0.85] at (1,-0.25) {lamp}; \pic[scale=0.85] at (2,-0.25) {nojuice};
\node[right,font=\scriptsize] at (2.36,0.1) {落ち込み};
\node[right,font=\scriptsize] at (2.13,3.75) {驚き！};
\node[left,font=\scriptsize] at (0.97,2.25) {ランプへのバースト};
\node[above,font=\scriptsize] at (2,1.75) {応答なし};
\draw[->,gray!70!black] (0,-0.75) -- (3.05,-0.75) node[right,black,font=\small] {$t$};
\draw[gray!60!black] (1,-0.8) -- (1,-0.7) node[below=2pt,font=\scriptsize] {ランプ、$1\,\text{s}$};
\draw[gray!60!black] (2,-0.8) -- (2,-0.7) node[below=2pt,font=\scriptsize] {ジュース、$2\,\text{s}$};
\end{tikzpicture}
\caption{三つの実験における\nt{DOPA}ニューロンの発火率（模式図、\cite{Schultz1997}による）。点線は一定の背景発火率。ランプは合図、しずくはジュース、斜線を引いたしずくは約束されたのに与えられなかったジュースである。応答は予測誤差~\eqref{eq:rpe}である。}
\label{fig:rpe}
\end{figure}

強化学習~\cite{SuttonBarto2018}を知っているプログラマなら、この量にすぐ気づくだろう。行動の選択を学習するエージェントは推定値$V(s)$を保持している。状態$s$にいるときに期待する報酬の量である。エージェントは各ステップの後、期待と実際に得たものを比べる。
\begin{empheq}[box=\keybox]{equation}
\delta_t \;=\; r_t \;+\; \gamma\,V(s_{t+1}) \;-\; V(s_t) ,
\label{eq:rpe}
\end{empheq}
そして推定値を修正する：$V(s_t)\leftarrow V(s_t)+\alpha\,\delta_t$。ここで$r_t$は得られた報酬、$\gamma<1$は割引率（将来の報酬が即時の報酬よりどれだけ安いか）、$\alpha$は学習率である。量$\delta_t$は\emph{時間差分誤差}と呼ばれる。

この量はドーパミンニューロンとまったく同じようにふるまう。予告なしのジュース：$V$はまだゼロで、$r>0$、$\delta>0$。学習済みのジュース：報酬はランプが予測しており、ジュースの直前の$V(s_t)$はすでに$r$に等しいので、$\delta=0$。ジュースが来ない：$V(s_t)=r$だが$r_t=0$なので、$\delta<0$。そしてバーストがランプに移るのは、まさにランプの瞬間に期待が跳ね上がるからである：$\gamma V(s_{t+1})-V(s_t)>0$。ここでのステップ$t,t+1,\dots$はアルゴリズム上の約束事にすぎない。生体の中に時間の刻みはなく、同じ量を連続時間で第\ref{sec:dofa}章で導く。

つまりドーパミンは、それをもとに学習すべきすべての相手に、スカラーの誤差信号$\delta$をブロードキャストで配信するものである。第一近似では、脳全体に1本の配線、各瞬間に一つの数である。この配線が実際にはどの程度ハーネスなのかは、第\ref{sec:dofaalt}章で検討する。

\section{その他のブロードキャストチャネル}

ほかの修飾物質については、その役割を同じ大域パラメータの言葉に翻訳する作業仮説があるだけである~\cite{Doya2002}。あくまで仮説として受け止めてほしい。

\emph{\nt{NORA}、ノルアドレナリン：ゲイン。}ノルアドレナリンニューロンはドーパミンニューロンよりさらに少なく、粗い推定で数万個だが、その出力は脳のほぼ全体に広がる。予想外のことや重要なことが起きると、これらは急激に活性化する。ノルアドレナリンは受け手のニューロンの伝達特性の傾きを変える。弱い入力はより弱く、強い入力はより強くなる。これは次のように測定される。受け手の背景スパイクが入力への応答より強く抑えられ、信号対雑音比が上がる~\cite{Foote1975NE}。単純なモデルではこれを一つの数で表す。ニューロンが入力電流$u$に発火率$f=F\bigl(g\cdot u\bigr)$で応答するなら、ノルアドレナリンは係数$g$を大きくする~\cite{AstonJones2005}（詳しくは第\ref{sec:nora}章）。$g$が大きいネットワークは「コントラスト高く」動作し、速く決定を下す。$g$が小さいと「ぼやけて」動作し、探索モードにある強化学習エージェントのように、より多くの選択肢を試す。

\emph{\nt{ACET}、アセチルコリン：学習率。}アセチルコリンは、ネットワークが現在の入力にどれだけ耳を傾け、自分が保存しているデータにどれだけ耳を傾けるかを制御する。アセチルコリンの濃度が高いと、感覚器からの入力が強まり、内部の「想起」結合が弱まる。同時に重みの変更が起こりやすくなる~\cite{Hasselmo2006}。大まかにいえば、これは「書き込み/読み出し」の切り替えスイッチであり、それとともに学習率$\alpha$でもある。

\emph{\nt{SERO}、セロトニン：計画の時間的視野。}セロトニンが何を伝えるかは最もよくわかっていない。仮説の一つは、これを式~\eqref{eq:rpe}の割引率$\gamma$と結びつける。セロトニンの濃度が高いことは、小さな報酬にすぐ飛びつかずに大きな報酬を待つ構えに対応する。セロトニンニューロンは一様にゆっくりと動作し、覚醒時には毎秒数スパイクで発火し、睡眠のある相ではほとんど黙り込む~\cite{JacobsAzmitia1992}。

6種類の神経伝達物質すべてを表~\ref{tab:transmitters}にまとめた。

\begin{table}[t]
\centering
\small
\begin{tabular}{>{\raggedright}p{2.4cm}>{\raggedright}p{3.0cm}>{\raggedright}p{2.0cm}>{\raggedright\arraybackslash}p{5.2cm}}
\toprule
ニューロンの型 & ニューロンの割合 & バス & 計算での役割 \\
\midrule
\nt{GLUT}（グルタミン酸） & 大多数 & アドレス & 正の重み、$w>0$ \\
\nt{GABA}（GABA） & 皮質で20--25\% & アドレス & 負の重み、$w<0$。安定化 \\
\nt{DOPA}（ドーパミン） & $\sim 6\cdot10^{-6}$ & ブロード\newline キャスト & 予測誤差$\delta$ \\
\nt{NORA}（ノルアドレナリン） & $\sim 6\cdot10^{-7}$ & ブロード\newline キャスト & ゲイン$g$（仮説） \\
\nt{ACET}（アセチルコリン） & 少数 & 両方 & 学習率$\alpha$。「書き込み/読み出し」（仮説） \\
\nt{SERO}（セロトニン） & $\sim 3\cdot10^{-6}$ & ブロード\newline キャスト & 割引$\gamma$（仮説） \\
\bottomrule
\end{tabular}
\caption{脳の主要な神経伝達物質を計算の言葉で表したもの。右の列は大幅な単純化である。グルタミン酸、GABA、ドーパミンについては信頼できるが、ほかは仮説である。}
\label{tab:transmitters}
\end{table}

\section{符号を決めるのは受け手}
\label{sec:receiver}

グルタミン酸は正の重み、GABAは負の重み、と言ったとき、実は少しごまかしていた。符号をもつ分子は一つもない。分子は単なる鍵である。それが捕らえられたとき何が起きるかを決めるのは\emph{錠}、つまり受け手の細胞にある受容体である。

最もよい例はドーパミンである。ドーパミンには二つのファミリーの受容体がある~\cite{Beaulieu2011}。一方では細胞の興奮を促し、他方では妨げる。行動を選択する領域では、一部のニューロンが第一の型の受容体を、別のニューロンが第二の型の受容体をもち、同じドーパミン信号が一方の群を強めると同時に他方の群を弱める。一つのブロードキャストパケットを、ネットワークのノードごとに異なって解釈するようなものである。

\begin{keyblock}
\begin{proposition}[メッセージの意味を決めるのは受け手]
\label{prop:receptor}
神経伝達物質の作用の符号と速さは、神経伝達物質そのものではなく、それが結合する受容体で決まる。受容体には二つの種類がある。\emph{イオンチャネル型}は、それ自体がイオンチャネルであり、1ミリ秒未満で開く。トランジスタのゲートのような、電流の直接制御である。\emph{代謝型}は細胞内で化学反応の連鎖を起こし、数百ミリ秒以上かけて作用する。むしろ設定レジスタへの書き込みに近く、細胞のその後の動作を変える。アドレスバスはおもに前者を通じて、ブロードキャストはおもに後者を通じて話す。
\end{proposition}
\end{keyblock}

それなら、なぜそもそも「グルタミン酸は興奮させる」と言えるのか。実際に出会うグルタミン酸受容体のほとんどすべてが興奮性であり、GABA受容体のほとんどすべてが抑制性だからである。これは自然法則ではなく、非常に信頼性の高い取り決めであり、規則~\eqref{eq:dalesign}はまさにその上に成り立っている。

\section{持ち帰るもの}

圧倒的多数のニューロンはアドレス指定のデータバスであり、重みの符号はニューロンごとに一つである。ごく少数派がブロードキャストチャネルであり、脳の広い領域にいくつかの共通の制御信号を送る。10万個のうちわずか2個ほどのニューロンに、残りのネットワーク全体がどう学習し、どのモードで動作するかがかかっている。エンジニアにとって、これは驚くことではない。どんなコンピュータでも制御レジスタはメモリセルよりも比べものにならないほど少ないが、それがなければマシンは動かない。

次章では、最も数の多い型である\nt{GLUT}ニューロンだけでできたネットワークが何を計算できるかを調べる。

\section{問題}

\begin{exercise}[\lvlA]
\label{ex:01:1}
\emph{山と配線。}表~\ref{tab:types}から（範囲は対数目盛での中央、\nt{ACET}の出力数は$10^5$）：(a)~ニューロン1個を人間1人とすると、\nt{GLUT}ニューロンは地球の人口の何倍になり、ブロードキャスト型全体はどれほどの規模の都市になるか。(b)~\nt{DOPA}、\nt{SERO}、\nt{NORA}、\nt{ACET}の終末の割合は、そのニューロンの割合の何倍か。
\end{exercise}

\begin{exercise}[\lvlA]
\label{ex:01:2}
\emph{ゼロへの復帰。}隙間のグルタミン酸は$c_0e^{-t/\tau_c}$で減衰する。$c_0=1\mM$、$\tau_c=1\ms$で、受け手は$10$~\textmu M以上の濃度を検知する。(a)~1回の放出の後、線路はどれだけの時間「使用中」になるか。(b)~ポンプが部分的に阻害され$\tau_c=10\ms$になった。線路が空くのが間に合う放出頻度の上限はいくらか。
\end{exercise}

\begin{exercise}[\lvlB]
\label{ex:01:3}
\emph{バーストの移動先。}図~\ref{fig:rpe}の実験で、ランプの後に状態$s_0\to\dots\to s_{K-1}$が刻み$\Delta$で続き、$T=K\Delta$、報酬$r$は最後の遷移で届き、ランプの時刻は予測できない。学習後の$V(s_k)$とランプへの応答$\delta$を求めよ。$\gamma=e^{-\Delta/\tau_\gamma}$とおき、$T=1$~s、$\Delta=0.1$~s、$\gamma=0.95$のときの$\tau_\gamma$を求めよ。
\end{exercise}

\begin{exercise}[\lvlB]
\label{ex:01:4}
\emph{シミュレーション：誤差による学習。}問題~\ref{ex:01:3}の連鎖を$K=10$、$\gamma=0.95$、$r=1$、規則$V(s_t)\leftarrow V(s_t)+\alpha\delta_t$でシミュレーションせよ。$\alpha=0.05$、$0.1$、$0.2$、$0.5$のとき、ランプへの応答がジュースへの応答を初めて上回る試行$n^*$はいくつか。$n^*$は$\alpha$にどう依存するか。
\end{exercise}

\begin{exercise}[\lvlB]
\label{ex:01:5}
\emph{設計：一つの数の配信。}数$\delta$を$10^{10}$個のニューロン全部に届けたい。中継役を含む各受け手は前の層から$10$個の終末を受け取る必要があり、1段ごとに$1\ms$かかる。ニューロン1個あたりの出力数が$10^4$（\nt{GLUT}）の場合と$10^{5.5}$（\nt{DOPA}）の場合で、最も経済的な中継ツリーのニューロン数と段数はいくつか。
\end{exercise}

\begin{exercise}[\lvlB]
\label{ex:01:6}
\emph{釣り合いはなぜ敏感か。}ネットワークは\nt{GLUT}が$80\,\%$、\nt{GABA}が$20\,\%$で、全対全に重み$J_E/N$と$-J_I/N$で結合し、$\tau\dot r=-r+Wr+h$である。$J_E=10$のとき、$W$のゼロでない唯一の固有値は$0.5$である。抑制性電流と興奮性電流の比を求め、ネットワークが雪崩に陥るには抑制を何パーセント弱めれば足りるかを求めよ。
\end{exercise}

\begin{exercise}[\lvlC]
\label{ex:01:7}
\emph{研究：\nt{SERO}は$\gamma$か。}問題~\ref{ex:01:3}より、ランプへの\nt{DOPA}ニューロンの応答は$re^{-T/\tau_\gamma}$である。遅延$T=1,\dots,5$~sをそれぞれ$n$回ずつ提示し、$\ln\delta$はばらつき$0.5$で測定される。$\tau_\gamma=2$~sと$2.4$~sを区別するには$n$はいくつ必要か（有意水準$0.05$、検出力$0.8$）。$\gamma$以外のどんな機構が$T$に対して同じ減衰を生むか。
\end{exercise}
